\uuid{MeXC}
\exo7id{4482}
\titre{exo7 4482}
\auteur{quercia}
\organisation{exo7}
\datecreate{2010-03-14}
\isIndication{false}
\isCorrection{true}
\chapitre{Série numérique}
\sousChapitre{Autre}
\cpge{mpsi-mp2i.nombres-reels-suites.d, pcsi.nombres-reels-suites.c, ptsi.nombres-reels-suites.c, mpsi-mp2i.procedes-sommatoires-discrets.a, pcsi.series-numeriques.a, ptsi.series-numeriques.a, bcpst1.suites-reelles, mp-mpi.topologie-espaces-normes.b, pc.espaces-vectoriels-normes.b, psi.espaces-vectoriels-normes.b, bcpst2.series-reelles}
\module{Analyse}
\niveau{L1}
\difficulte{}

\contenu{
\question{
Soit $(u_n)$ une suite de complexes qui converge au sens
de Césaro vers zéro.

Étudiez la suite de terme général
$v_n=\sum_{k=0}^n \frac{u_k}{n+k+1}\cdotp$
}
\reponse{
Transformation d'Abel $+$ découpage, $v_n\to0$ lorsque $n\to\infty$.
}
}
